\section{Vektorgeometrie}

Ein \emph{Vektor} $\vec{v}$ beschreibt eine gerichtete Translation im Raum und ist durch Richtung und Betrag $|\vec{v}|$ eindeutig bestimmt. Besondere Vektoren:
\begin{itemize}
     \item \textbf{Nullvektor} \(\vec{0}\): Ein Vektor mit Länge 0, der keine Richtung hat.
     \item \textbf{Einheitsvektor} \(\vec{e}_i\): Ein Vektor, der in der Richtung der \(i\)-ten Achse liegt und die Länge 1 hat.
     \item \textbf{Gegenvektor} \(-\vec{v}\): Ein Vektor, der die gleiche Länge wie \(\vec{v}\) hat, aber in die entgegengesetzte Richtung zeigt.
\end{itemize}


\subsection{Rechenregeln}

Vektoren lassen sich addieren und mit Skalaren $\lambda \in \mathbb{R}$ multiplizieren; beide Operationen sind komponentenweise definiert.

\textbf{Addition:}
\[
     \vec a + \vec b = \begin{pmatrix} a_1 + b_1 \\ \vdots \\ a_n + b_n \end{pmatrix}
\]

\textbf{Skalare Multiplikation}
\[
     \lambda \vec a = \begin{pmatrix} \lambda a_1 \\ \vdots \\ \lambda a_n \end{pmatrix}
\]

\subsection{Linearkombination}
Eine \emph{Linearkombination} von Vektoren $\vec a_1, \ldots, \vec a_n$ ist ein Ausdruck der Form
\[
     \vec v = \sum_{i=1}^{n} \lambda_i \vec a_i, \quad \lambda_i \in \mathbb{R}
\]

\subsection{Lineare Unabhängigkeit}
Vektoren $\vec{a}_1, \ldots, \vec{a}_n$ sind \emph{linear unabhängig}, wenn:
\[
     \vec{0} = \sum_{i=1}^{n} \lambda_i \vec{a}_i \quad \Rightarrow \quad \lambda_1 = \lambda_2 = \cdots = \lambda_n = 0
\]

\subsection{Kollinearität}
Zwei Vektoren $\vec a$ und $\vec b$ sind \emph{kollinear} (parallel), wenn einer ein Vielfaches des anderen ist:
\[
     \vec a \parallel \vec b \;\Leftrightarrow\; \exists\,\lambda\colon \vec a = \lambda\vec b
\]

\subsection{Komplanarität}
Vektoren liegen in derselben Ebene, wenn sie als Linearkombination von zwei nicht-kollinearen Vektoren dargestellt werden können:
\[
     \vec c = \lambda\vec a + \mu\vec b, \quad \text{mit } \vec a \nparallel \vec b
\]

\subsection{Koordinaten und Betrag}

Jeder Vektor im $\mathbb{R}^n$ wird durch $n$ reelle Komponenten bezüglich der Standardbasis dargestellt; sein Betrag ergibt sich aus dem Satz des Pythagoras.

\[
     \vec a = \sum_{i=1}^{n} a_i \vec e_i =
     \begin{pmatrix} a_1 \\ \vdots \\ a_n \end{pmatrix}, \qquad
     |\vec a| = \sqrt{\sum_{i=1}^{n} a_i^2}
\]
Der Verbindungsvektor von Punkt $P$ zu Punkt $Q$ und dessen Länge ist gegeben durch:
\[
     \overrightarrow{PQ} = \vec r(Q) - \vec r(P), \qquad
     \overline{PQ} = |\overrightarrow{PQ}|
\]



\subsection{Skalarprodukt}

Das \emph{Skalarprodukt} zweier Vektoren liefert eine reelle Zahl und misst den "Anteil" von $\vec{b}$ in Richtung $\vec{a}$. Es ist kommutativ und distributiv.

\[
     \vec a \cdot \vec b = \sum_{i=1}^{n} a_i b_i = |\vec a|\,|\vec b|\cos\varphi, \qquad 0 \le \varphi \le \pi
\]

\begin{center}
\begin{tikzpicture}[scale=0.9,>=latex]
     % vectors
     \coordinate (O) at (0,0);
     \coordinate (A) at (3,0);
     \coordinate (B) at (1.6,1.7);
     % projection of b onto a (foot point)
     \coordinate (P) at (1.6,0);

     % projection length helper line
     \draw[dashed,gray] (B) -- (P);

     % vector a
     \draw[->,thick,blue] (O) -- (A) node[below right] {$\vec a$};
     % vector b
     \draw[->,thick,red]  (O) -- (B) node[above] {$\vec b$};

     % projection b_a on a
     \draw[->,very thick,red!60!black] (O) -- (P)
          node[midway,below] {$\vec b_a$};

     % angle phi
     \pic[draw,->,"$\varphi$",angle radius=8mm,angle eccentricity=1.35]
          {angle=A--O--B};

     % right angle at foot point
     \draw (P) ++(-0.18,0) -- ++(0,0.18) -- ++(0.18,0);
\end{tikzpicture}
\end{center}

\textbf{Eigenschaften:}
\begin{itemize}
     \item Orthogonalitätstest: $\vec a \cdot \vec b = 0 \;\Leftrightarrow\; \vec a \perp \vec b$
     \item $\vec a \cdot \vec a = |\vec a|^2$
\end{itemize}


\textbf{Winkel:}
\[
     \varphi = \arccos\!\left(\frac{\vec a \cdot \vec b}{|\vec a|\,|\vec b|}\right)
\]

\textbf{Orthogonale Projektion} von $\vec b$ auf $\vec a$:
\[
     \vec b_a = \frac{\vec a \cdot \vec b}{|\vec a|^2}\,\vec a, \qquad
     |\vec b_a| = \frac{|\vec a \cdot \vec b|}{|\vec a|}
\]

\subsection{Vektorprodukt ($\mathbb{R}^3$)}

Das \emph{Vektorprodukt} zweier Vektoren im $\mathbb{R}^3$ liefert einen neuen Vektor, der senkrecht auf beiden steht. Die Richtung folgt der Rechte-Hand-Regel.

\[
     \vec a \times \vec b =
     \begin{pmatrix} a_2b_3-a_3b_2 \\ a_3b_1-a_1b_3 \\ a_1b_2-a_2b_1 \end{pmatrix}, \qquad
     |\vec a \times \vec b| = |\vec a|\,|\vec b|\sin\varphi
\]

\textbf{Eigenschaften:}
\begin{itemize}
     \item $\vec a \times \vec b \perp \vec a, \vec b$ (Richtung via Rechte-Hand-Regel)
     \item $|\vec a \times \vec b|$ = Fläche des aufgespannten Parallelogramms
     \item $\vec a \times \vec b = \vec 0 \;\Leftrightarrow\; \vec a \parallel \vec b$
     \item $\vec a \times \vec b = -(\vec b \times \vec a)$ (Antikommutativität)
     \item $\vec a \times (\vec b + \vec c) = \vec a \times \vec b + \vec a \times \vec c$ (Distributivität)
     \item $\lambda(\vec a \times \vec b) = (\lambda\vec a) \times \vec b = \vec a \times (\lambda\vec b)$ (Homogenität)
     \item $\vec e_1 \times \vec e_2 = \vec e_3$, \quad $\vec e_2 \times \vec e_3 = \vec e_1$, \quad $\vec e_3 \times \vec e_1 = \vec e_2$
\end{itemize}